\scanpage{42}
Soit $S$ un ens fini\inserted{Cardinal\unsure{,} $N$}, avec une structure polygonale ou \unsure{bi-circularisation} $\{u,u^{-1}\}$ ($\mathrm{Card}\,S\geq\unsure{3|2}$).\unsure{,}

\deleted{\uline{\unsure{Ax.|Déf.}}} $T\subset S$ \inserted{de cardinal $i$} est dit un \deleted{\unsure{ens}} \uline{segment} de $S$, si

a) \deleted{\unsure{non|on}} $T\ne S$, $T\ne\varnothing$ \deleted{$i<N$} i.e. $0<i<N$

b) $\exists\,s_0\in S$, \deleted{$i\in\mathbb N$}, tel que $T=\{s_0,us_0,\ldots,u^{i-1}s_0\}$

(donc, posant $t_0={\unsure{w|u}}^{i-1}s_0$, on aura aussi $T=\{t_0,vt_0,\ldots,v^{i-1}t_0\}$, où $v=u^{-1}$).

\uline{Ex.} les parties réduites \unsure{:|à} un point.

les \unsure{parties-arêtes}

% Slanted, cancelled proposition in the left margin.
\margin{\deleted{\uline{\unsure{Propr.|Prop.}} \unsure{Soit|Soient} $A\unsure{,}A'$ deux \unsure{segments},}}

$T$ est un segm. \unsure{sss|ssi} $S\setminus T$ en est un.

\inserted{\unsure{\uline{Intersection}|Intersection} de deux segments \inserted{\unsure{\uline{contenus dans un segment}|contenus dans un segment}} en est un sss il est $\ne\varnothing$}

\inserted{Réunion de deux segments \unsure{pris ds un contexte|pris ds un contre-ex.|qui se rencontrent}, et tels que $T'\cup T''\ne S$ \inserted{est un segment}}
% Slanted annotation at the right edge, attached near the inserted restriction.
\margin{i.e de \deleted{\unsure{réunion}} $\ne S$}

Extrémités d'un segment $s_0,t_0$ (distinctes sss $i\ne1$ i.e. $T$ non réduit à un point).

Origine, extrémité pour une . orientation.

[ \uline{NB} Si $u$ i.e. l'orientation choisi, le $s_0$ est unique i.e. il y a corr. 1-1 entre l'ens des segments de $S$, et l'ens
\[
 (S,[1,N-1]_{\mathbb N}),\ \text{via}\quad (s_0,i)\longmapsto\{s_0,us_0,\ldots u^{i-1}s_0\}.
\]

Soit $S'\subsetneq S$ une partie de $S$. Alors \deleted{les segments de $S$ contenus} deux segments de $S'$ contenus ds $S'$, et qui se rencontrent, ont comme réunion un segment \inserted{de $S$} contenu dans $S'$ -- \unsure{cel|cela} implique
% archive sheet number: 41
